\subsection{BCP: impacto de negocio y niveles de servicio}\label{sec:sd8-bcp-bia}
El BCP aplica el ciclo de ISO 22301: identificar procesos y dependencias, analizar impacto, seleccionar continuidad, ejecutar una respuesta autorizada, comprobar recuperación y mejorar con resultados. El DRP aplica ISO/IEC 27031 a los recursos TIC de esa misma cadena. No se declara certificación ni simulacro realizado. Proyecto conserva versiones; SRE mantiene procedimientos y Seguridad custodia accesos; la Contraparte valida condiciones funcionales y prioridades. \parencites[Formulario T-7, apartado 10; arts.~20 y 78]{sd8-ba}[RT-10.02--04]{sd8-bt}

La Tabla~\ref{tab:sd8-bia} clasifica por consecuencia, no por tecnología. Las funciones críticas incluyen captura y consulta que permiten demostrar custodia y seguridad; una salida administrativa o un informe diferible puede pertenecer a otra clase dentro del mismo dominio. Para servicios críticos se comprometen RTO máximo de cuatro horas y RPO regional máximo de quince minutos. En captura desconectada, el requisito es conservar íntegramente los hechos críticos: ese RPO no permite descartar quince minutos de hechos locales. Las autonomías del caso son límites de capacidad de operación, no autorización para continuar una maniobra insegura.
\begingroup\setlength{\tabcolsep}{3pt}
\begin{longtable}{L{.22\textwidth}L{.10\textwidth}L{.29\textwidth}L{.29\textwidth}}
\caption{Impacto de negocio, dependencia y modo seguro}\label{tab:sd8-bia}\\\hline
\EncabezadoTabla{Servicio} & \EncabezadoTabla{Clase} & \EncabezadoTabla{Impacto/dependencia} & \EncabezadoTabla{Continuidad y límite}\\\hline\endfirsthead
\hline\EncabezadoTabla{Servicio} & \EncabezadoTabla{Clase} & \EncabezadoTabla{Impacto/dependencia} & \EncabezadoTabla{Continuidad y límite}\\\hline\endhead
Custodia, amarras, salida/regreso y legajo de zarpe & Crítica & Entrega indebida o desconocer quién está en el agua; autoridad local, diario y radio. & Recinto 24 h sin enlace; legajo local y coordinación humana; sin autorización no se concede zarpe ni entrega.\\\hline
Escuela: presencia, retirada, salud y emergencias & Crítica & Menor sin ubicación o entrega no autorizada; lista vigente, permisos y monitor. & Recinto 24 h y terminal 12 h; comprobación humana inmediata; si no hay certeza se suspende nueva salida y se aplica emergencia.\\\hline
Despacho de combustible y seguridad de maniobra & Crítica & Daño físico y pérdida de trazabilidad; habilitación, control local y condiciones seguras. & Captura local 24 h; detener despacho/izaje si falla la condición de seguridad o el registro exigible.\\\hline
Aviso urgente y contacto de afectados & Crítica & Falta de instrucción a personas expuestas; destinatarios, acuse, radio/telefonía alternativa. & Registro de llamada por no acuse; no presumir entrega ni cobertura marítima continua.\\\hline
Gestión de faena, acreditación y entrega de residuos & Alta & Se impide faena o entrega sin habilitación/trazabilidad; manifiesto, gestor y permisos. & Captura de terreno 12 h; detener operación dependiente sin autorización; conservar entrega/recepción probatorias.\\\hline
Administración: cierre, cobros y contabilidad & Alta & Diferencia de saldo o falta de expediente; legado, suplente funcional e interfaz externa. & Consulta local del último estado con antigüedad visible; diferir efectos externos y conciliar; no inferir saldo vigente.\\\hline
Indicadores ambientales y analítica & Media & Evidencia o beneficio incompletos; serie, asociación y corte verificable. & Conservar captura prioritaria; diferir cálculo/exportación y documentar atraso sin inventar resultado.\\\hline
Contenido institucional no transaccional & Baja & Información ordinaria diferida; no comprende portal de atención o avisos urgentes. & Mensaje de indisponibilidad y recuperación posterior; sin afectar canales críticos.\\\hline
\end{longtable}\endgroup
\FuenteTabla{Clasificación propuesta de Synaptix; Caso 07, caps.~3/15/19, y BA art.~78.}
Cada clase mantiene disponibilidad mensual mínima, respuesta y resolución: crítica 99,9\,\%, 15 min y 4 h; alta 99,5\,\%, 1 h y 8 h; media 99,0\,\%, 4 h y 24 h; baja 98,0\,\%, 8 h y 48 h. Crítica/alta disponen de atención 24\,$\times$\,7; media, horario hábil extendido; baja, horario hábil. Un incidente crítico conserva informe de causa raíz en cinco días hábiles y uno alto en diez. Disponibilidad se mide sobre la transacción de extremo a extremo; éxito de nube, réplica o pantalla no sustituye recepción y efecto de negocio. \parencite[art.~78]{sd8-ba}

\subsection{Desconexión, falla local y captura alternativa}\label{sec:sd8-bcp-local}
\textbf{BC-01, pérdida WAN (R-02/R-03/R-08).} SRE distingue enlace de energía y nodo. Se retiran los dos enlaces en el ensayo; la Célula Local mantiene recepción, asignación, salida/regreso, legajo, combustible y escuela durante 24 horas, con permisos locales autorizados y diario durable. Pantalán, varadero y escuela conservan 12 horas en terminales; Ilque, ocho horas sin cobertura. Se advierte antigüedad de copias y funciones externas pendientes. No se suponen 24 horas de batería de cada dispositivo ni combustible de generador por perder Internet. Energía y redundancia mantienen el diseño dimensionado de SD4, sin ampliar su autonomía por esta declaración.

En mostrador se identifica embarcación/persona, se verifica documentación local vigente, se prepara legajo en no más de 90 segundos y se registra el hecho con ID, secuencia, operador, instante y resultado. Si la autoridad requiere un canal indisponible, el legajo queda preparado y su remisión se realiza por el medio autorizado disponible; la plataforma no otorga zarpe ni interpreta silencio como autorización. No están disponibles confirmación de pagos/saldos externos, consulta actual a una autoridad caída, tratamiento remoto de IN3, federación remota ni analítica central fresca; se difiere el efecto y se muestra su estado, se usa captura/revisión manual de póliza cuando proceda y se conserva la tarea. El permiso local vigente no se sustituye por una contraseña compartida ni por una autenticación remota inexistente. \parencites[RT-03.10/13; RT-09.01; numeral 16.1, decisión 21]{sd8-c07}[RT-03.13]{sd8-bt}

Al reconectar, Sincronización envía automáticamente hechos desde su diario, verifica huella/secuencia e idempotencia, aplica autoridad por dominio y concilia sin pérdida. La prueba exige que después de 24 h de desconexión termine en 30 minutos como máximo, sin intervención manual en esa sincronización y sin perder salida, regreso, despacho o asistencia. Un error conserva el original y abre incidente; no se resuelve con última escritura que borre un hecho de seguridad, ni se comunica éxito porque la cola esté vacía. Un registro alternativo en papel es otro escenario y exige incorporación controlada; no cuenta como prueba de operación electrónica desconectada.

\textbf{BC-02, falla de nodo local (R-02/R-13).} SRE detecta mediante salud de servicio y transacción, aísla el nodo que falla y promueve únicamente la réplica local que conserve corte y exclusión del escritor anterior. Verifica energía, red, diario, permisos y captura completa antes de declarar servicio. Si falla también la alternativa o no puede excluir escritura doble, se bloquea la función afectada, se preserva almacenamiento y se activa respaldo manual seguro; no se confunde esa condición con simple pérdida WAN ni se exige administración técnica a la marina. Reparación y resiembra son tareas Synaptix; el titular de proceso valida negocio.

\textbf{BC-03, incertidumbre sobre alumnos, custodia o navegación (R-10/R-11).} El referente de escuela y su suplente contrastan lista local, salida real, acompañante, autorización y retirada. Ante discrepancia se mantiene vigilancia y se activa la respuesta humana del proceso; no se inventa presencia a partir de matrícula ni ausencia de mensaje. Operaciones contrasta salida/regreso y custodia por radio y contacto, sin concluir emergencia ni regreso automáticamente por telemetría tardía. Sin autorización probatoria no se entrega un menor o embarcación. La comunicación no exige tocar pantallas durante amarre o izaje; se realiza antes/después en condición segura.

\textbf{BC-04, aviso masivo fallido o marejada (R-01/R-13).} Notificaciones conserva destinatario, intento, acuse técnico y confirmación humana separados. Sin acuse, el operador abre llamada con titular/suplente, instrucción, hora, respuesta y siguiente escalamiento. Radio y teléfono disponible sirven de alternativa para expuestos, con evidencia de intento y contacto efectivo. Aviso oficial de marejada suspende toda actividad programada del proyecto; la respuesta a una emergencia real conserva su procedimiento. Se comunica al CLIENTE afectación, medidas y próxima actualización; no se utiliza continuidad para seguir una faena prohibida.

\textbf{BC-05, ausencia de conocedor del legado o cierre fallido (R-04/R-12).} Datos usa copia/inventario y procedimiento de cierre transferidos desde mes 6, con suplente técnico; Administración usa el suplente funcional ensayado. Se conserva último corte comprobado, se verifica recuento/saldo y se mantiene convivencia mientras no haya evidencia suficiente. No se calcula deuda por inferencia ni se reemplaza interpretación con un respaldo sin probar. La conciliación diaria y cierre mensual durante marcha blanca mantienen las puertas de SD5; un dato no explicado bloquea el corte.

Todo registro alternativo usa folio único, proceso, hecho, operador, hora, causa y relación con autorización; contiene sólo los datos imprescindibles. El titular funcional custodia físicamente la hoja bajo acceso restringido y SRE/Datos incorporan una imagen y el hecho trazable al ID original, con revisión independiente de recuento, efectos y duplicados. Se separa captura de validación: incorporar un papel no permite crear un segundo cargo, salida o despacho. El retorno requiere conciliar todos los folios, resolver pendientes y conservar evidencia por la retención del dominio de SD5; no se reinstala un cuaderno como sistema oficial permanente. La contingencia manual reduce daño pero no acredita la disponibilidad electrónica ni subsana una brecha del diseño.
